\problemname{Neuralt Bakverk}
\noindent

Ett neuralt bakverk är en riktad graf med $N$ noder, vars funktion är att baka kakor. 
I början läggs $X$ enheter ingredienser på nod $1$, medan de andra noderna får noll 
ingredienser. Därefter sker följande varje sekund för varje nod:


\begin{itemize}
\item Låt $a_v$ vara antalet ingredienser på nod $v$, och låt $d_v$ vara utgraden hos nod $v$.
\item Om $d_v = 0$ så händer ingenting. Annars, skriv $a_v = d_v\cdot q + r$ där $0 \leq r < d_v$. Skicka sedan $q$ enheter ingredienser till var och en av noderna som $v$ pekar på, och behåll resterande $r$.
\end{itemize}

Med andra ord så delar varje nod upp ingredienserna lika så gått det går och skickar 
vidare. Varje sekund så sker allt detta samtidigt för alla noder, så det spelar 
ingen roll i vilken ordning som noderna betraktas.

De ingredienser som hamnar på nod $N$ är de färdiga kakorna. 
För att vara ett giltigt neuralt bakverk så måste $d_N = 0$. 
En nod får inte heller peka på sig själv, eller peka på samma nod mer än en gång.

Du får givet heltalen $X$ och $Y$, och din uppgift är att bygga ett neuralt 
bakverk sådant att exakt $Y$ kakor bakas om $X$ enheter ingredienser används. 
Mer specifikt så kommer bakverket att köras i $10^4$ steg, och om nod $N$ då 
har $Y$ kakor så har du klarat testfallet. Du får poäng beroende på hur litet $N$ är.


\section*{Indata}

Indatan består av en rad med två heltal $X$ och $Y$ ($1 \leq X,Y \leq 10^{18}$).

\section*{Utdata}

Om det inte finns något neuralt bakverk som bakar $Y$ kakor med $X$ ingredienser, skriv 
ut $-1$.

Annars, skriv först ut en rad med två heltal $N,M$, antalet noder och kanter i ditt nätverk. 
Skriv därefter ut $M$ rader som var och en innehåller två heltal $u_i, v_i$ ($1 \leq u_i, v_i \leq N$),
som indikerar att en kant går från nod $u_i$ till nod $v_i$.

\section*{Poängsättning}

På var och en av de två testgrupperna kommer du få poäng enligt formeln

$$TODO$$

där $g$ är testgruppens poäng och $N$ är hur många noder ditt nätverk har.

\noindent
\begin{tabular}{| l | l | p{12cm} |}
  \hline
  \textbf{Grupp} & \textbf{Poäng} & \textbf{Gränser} \\ \hline
  $1$    & $20$        & $X,Y \leq 20$ \\ \hline
  $2$    & $80$        & Inga ytterligare begränsningar. \\ \hline
\end{tabular}

\section*{Förklaring av exempelfall}

